% experimental-setup.tex -- Enterprise Node Orchestrator Paper (RU-V5)

\section{Experimental Setup}

\subsection{Benchmark Methodology}

The experimental evaluation consists of two complementary benchmarks:
(1)~an \emph{orchestration overhead benchmark} that measures per-phase
latencies across the complete proof cycle, and (2)~a \emph{pipeline
comparison benchmark} that quantifies the throughput improvement of
pipelined over sequential execution.

\subsubsection{Orchestration Overhead Benchmark}

Five scenarios are evaluated, each varying the batch size and prover
backend while holding the circuit depth constant at 32 levels:

\begin{enumerate}
\item \texttt{orchestration\_overhead\_only}: Batch 8, 0~ms proving
  (isolates pure orchestration cost). 50 iterations.
\item \texttt{snarkjs\_d32\_b8}: Batch 8, 12,757~ms proving (snarkjs
  WASM, 274K constraints). 30 iterations.
\item \texttt{rapidsnark\_d32\_b8}: Batch 8, 2,500~ms proving
  (rapidsnark, 274K constraints). 30 iterations.
\item \texttt{snarkjs\_d32\_b16}: Batch 16, 28,000~ms proving (snarkjs
  WASM, $\sim$548K constraints). 30 iterations.
\item \texttt{rapidsnark\_d32\_b64}: Batch 64, 12,000~ms proving
  (rapidsnark, $\sim$2.2M constraints). 30 iterations.
\end{enumerate}

Each scenario executes 3 warm-up iterations (discarded) followed by the
measurement iterations. Per-phase timings are recorded for: batch
formation, witness generation, proving, DAC attestation, L1 submission,
and checkpoint. Orchestration overhead is computed as the sum of batch
formation and witness generation.

\subsubsection{Pipeline Comparison Benchmark}

Six configurations are simulated with 10 sequential batches each,
comparing sequential and pipelined architectures. Component latencies
are calibrated from prior research unit measurements:

\begin{itemize}
\item SMT insert: 1.8~ms/tx (RU-V1, Poseidon hash, depth 32)
\item WAL write: 0.149~ms/tx (RU-V4, group commit)
\item Batch formation: 0.02~ms (RU-V4)
\item Witness generation: 72~ms/tx (RU-V2, 578~ms for 8~txs)
\item DAC attestation: 163~ms (RU-V6, JavaScript, 500~KB)
\item L1 submission: 2,000~ms (estimated Avalanche Fuji block time)
\end{itemize}

Proving times are parameterized per backend: snarkjs (12,757~ms for b8,
28,000~ms for b16, 150,000~ms for b64) and rapidsnark (2,500~ms for b8,
5,000~ms for b16, 12,000~ms for b64).

\subsection{Mock Calibration}

Stage~1 benchmarks use mock components calibrated to measured latencies
from prior experiments to isolate orchestration overhead from component
variability. Each mock implements the same TypeScript interface as the
production component:

\begin{itemize}
\item \textbf{MockSMT}: Simulates Poseidon hash latency via CPU-bound
  SHA-256 iterations (32 rounds per insert) with spin-wait padding
  to match the 1.0--1.8~ms target from RU-V1.

\item \textbf{MockBatchBuilder}: Simulates witness generation at 72~ms/tx,
  scaling linearly with batch size. Calibrated to the 578~ms measurement
  from RU-V2 for depth-32, batch-8 configurations.

\item \textbf{MockProver}: Configurable timer-based delay returning a
  valid Groth16 proof structure. Default: 12,757~ms (snarkjs d32, b8).

\item \textbf{MockDAC}: 163~ms delay matching the JavaScript Shamir
  attestation measurement from RU-V6.

\item \textbf{MockL1Submitter}: 2,000~ms delay matching the estimated
  Avalanche Fuji block confirmation time.
\end{itemize}

\subsection{Statistical Protocol}

For each scenario, we report mean, standard deviation, minimum, maximum,
and P95 latency. The 95\% confidence interval (CI) half-width must be
within 10\% of the mean for the result to be accepted. All replications
use independent random seeds for transaction generation.

\subsection{Platform}

All benchmarks execute on Node.js v22 (Windows 11, x64) with commodity
hardware. Memory footprint is measured via \texttt{process.memoryUsage().heapUsed}.

\subsection{State Machine Validation}

A dedicated test suite validates the state machine implementation with
46 test cases covering: happy path transitions, pipelined receiving during
Proving and Submitting, error recovery from all states, invalid transition
rejection, transition history tracking, and event emitter correctness.
